\chapter{Free Electrical Oscillations: RLC}\label{ch:g12:rlc-oscillations}

Turn the heavy dial of an attic radio and the stations file past the
needle: a voice, violins, static, a voice again. Behind the panel there
is no \omterm{def:g11:work-of-force:sign}{motor} and no computer --- a \omterm{def:g11:magnetic-fields:solenoid}{coil}, a \omterm{def:g12:rc-rl-circuits:capacitor}{capacitor} whose interleaved
plates the dial rotates, and a length of aerial wire. This chapter
connects last chapter's two \omterm{def:g10:energy-conservation:energy}{energy} tanks (\cref{ch:g12:rc-rl-circuits})
head to head and finds the fastest pendulum ever built: charge swinging
back and forth, up to millions of times per second, at a tempo the dial
chooses.

\section{An electrical pendulum: the LC circuit}

Charge a \omterm{def:g12:rc-rl-circuits:capacitor}{capacitor} $C$ to a \omterm{def:g11:circuits-and-power:voltage}{voltage} $U_0$ --- its plates then hold
$\pm Q_0$ with $Q_0 = C U_0$ --- and at $t = 0$ switch it onto an
\emph{ideal} \omterm{def:g11:magnetic-fields:solenoid}{coil} $L$, of zero \omterm{def:g11:circuits-and-power:resistance}{resistance}. No source anywhere in the
loop: whatever happens next, the circuit does on its own.

\begin{omfigure}
\begin{circuitikz}[font=\small]
  \draw (0,0) to[C=$C$, v=$u$] (0,2.4) to[nos, l=$K$] (2.6,2.4)
    -- (3.8,2.4) to[L=$L$, i=$i$] (3.8,0) -- (0,0);
  \node[right, font=\footnotesize] at (0.18,1.5) {$+q$};
  \node[right, font=\footnotesize] at (0.18,0.9) {$-q$};
\end{circuitikz}

{\small The \omterm{prop:g12:rlc-oscillations:equation}{LC circuit}: a \omterm{def:g12:rc-rl-circuits:capacitor}{capacitor} charged to $Q_0$, a switch, an
ideal \omterm{def:g11:magnetic-fields:solenoid}{coil}. Closing $K$ lets the charge spill through the \omterm{def:g11:magnetic-fields:solenoid}{coil} --- and
the \omterm{def:g11:magnetic-fields:solenoid}{coil}, hating change, will not let it stop.}
\end{omfigure}

\begin{proposition}[The LC equation]\label{prop:g12:rlc-oscillations:equation}
\index{LC circuit}
In the ideal LC loop the charge $q(t)$ of the upper plate obeys
\[
L\,\frac{d^2q}{dt^2} + \frac{q}{C} = 0,
\qquad\text{i.e.}\qquad
\frac{d^2q}{dt^2} = -\frac{1}{LC}\,q .
\]
\end{proposition}

\begin{proof}
Around the loop the \omterm{def:g11:circuits-and-power:voltage}{voltages} sum to zero
(\cref{prop:g11:circuits-and-power:laws}): $u_L + u_C = 0$. The
\omterm{def:g12:rc-rl-circuits:capacitor}{capacitor} gives $u_C = q/C$ and the \omterm{def:g11:magnetic-fields:solenoid}{coil} $u_L = L\,di/dt$
(\cref{def:g12:rc-rl-circuits:inductor}), with $i = dq/dt$ counting
the charge arriving on the upper plate
(\cref{prop:g12:rc-rl-circuits:cdudt}); substitute.
\end{proof}

\begin{theorem}[Free oscillations of the LC circuit]\label{thm:g12:rlc-oscillations:oscillation}
\index{electrical oscillations}
Released at $t = 0$ with $q(0) = Q_0$ and $i(0) = 0$, the circuit
oscillates:
\[
q(t) = Q_0 \cos\!\left(\frac{2\pi t}{T_0}\right),
\qquad
i(t) = -I_0 \sin\!\left(\frac{2\pi t}{T_0}\right),
\qquad
T_0 = 2\pi\sqrt{LC},
\]
with peak current $I_0 = 2\pi Q_0/T_0 = Q_0/\sqrt{LC}$.
\end{theorem}

\begin{proof}
Verify by substitution. Differentiating,
$i = dq/dt = -\frac{2\pi}{T_0} Q_0 \sin(2\pi t/T_0)$ --- the stated
$i(t)$ --- and again $\frac{d^2q}{dt^2} = -(\frac{2\pi}{T_0})^2 q$,
which equals $-q/(LC)$ exactly when $(2\pi/T_0)^2 = 1/(LC)$, i.e.\
$T_0 = 2\pi\sqrt{LC}$. And $q(0) = Q_0$, $i(0) = 0$. That no
\emph{other} function fits the equation and the start is admitted,
exactly as for the spring
(\cref{rem:g12:oscillators-and-time:uniqueness}).
\end{proof}

\begin{definition}[Natural period and frequency]\label{def:g12:rlc-oscillations:period}
The \emph{natural period}\index{natural period!of an LC circuit} of an
\omterm{prop:g12:rlc-oscillations:equation}{LC circuit} is $T_0 = 2\pi\sqrt{LC}$; its
\emph{natural frequency}\index{natural frequency!of an LC circuit} is
\[
f_0 = \frac{1}{T_0} = \frac{1}{2\pi\sqrt{LC}} .
\]
Dimensional check: $\qty{1}{H} = \qty{1}{V.s/A}$ and
$\qty{1}{F} = \qty{1}{C/V}$, so $LC$ carries
$\unit{s} \times \unit{C/A} = \unit{s^2}$ --- $\sqrt{LC}$ is a time,
as it must be.
\end{definition}

\begin{example}[First numbers]\label{ex:g12:rlc-oscillations:firstnumbers}
$L = \qty{10}{mH}$ with $C = \qty{1.0}{\micro F}$:
$\sqrt{LC} = \qty{1.0e-4}{s}$, so $T_0 = \qty{0.63}{ms}$ and
$f_0 \approx \qty{1.6}{kHz}$ --- an audible tone. $L = \qty{1.0}{mH}$
with $C = \qty{100}{pF}$: $T_0 = \qty{2.0}{\micro s}$,
$f_0 \approx \qty{0.50}{MHz}$ --- a radio \omterm{def:g12:oscillators-and-time:oscillator}{frequency}. Only the
\emph{product} $LC$ sets the tempo.
\end{example}

\begin{remark}[A quarter period out of step]\label{rem:g12:rlc-oscillations:quarter}
The current is a quarter \omterm{def:g12:oscillators-and-time:oscillator}{period} out of \omterm{def:g12:oscillators-and-time:phase}{phase} with the charge: $i = 0$
when $q = \pm Q_0$ (plates full, flow reversing) and $i = \pm I_0$
when $q = 0$ (plates empty, flow at full tilt) --- exactly the
pendulum, motionless at the extremes and fastest at the bottom.
\end{remark}

\begin{omfigure}
\begin{tikzpicture}
\begin{axis}[omaxis, width=9.2cm, height=5.0cm, xlabel={$t$},
    xmin=0, xmax=2, ymin=-1.4, ymax=1.75, xtick={0,0.5,1,1.5,2},
    xticklabels={$0$,$T_0/2$,$T_0$,$3T_0/2$,$2T_0$},
    ytick={-1,0,1}, legend pos=north east]
  \addplot[omDef, thick, domain=0:2, samples=200] {cos(360*x)};
  \addplot[omThm, thick, domain=0:2, samples=200] {-sin(360*x)};
  \legend{$q/Q_0$, $i/I_0$}
\end{axis}
\end{tikzpicture}

{\small Charge and current of the free \omterm{prop:g12:rlc-oscillations:equation}{LC circuit}: the current peaks
each time the charge crosses zero, a quarter \omterm{def:g12:oscillators-and-time:oscillator}{period} out of step.}
\end{omfigure}

\section{The energy waltz}

\begin{proposition}[Conservation of the stored energy]\label{prop:g12:rlc-oscillations:energy}
\index{energy!of an LC circuit}
In the ideal \omterm{prop:g12:rlc-oscillations:equation}{LC circuit} the total stored \omterm{def:g10:energy-conservation:energy}{energy}
(\cref{prop:g12:rc-rl-circuits:energy})
\[
E \;=\; E_C + E_L \;=\; \frac{q^2}{2C} + \tfrac12\,L i^2
\;=\; \frac{Q_0^2}{2C} \;=\; \tfrac12\,L I_0^2
\]
is constant: entirely in the \omterm{def:g12:rc-rl-circuits:capacitor}{capacitor} when $q = \pm Q_0$, entirely in
the \omterm{def:g11:magnetic-fields:solenoid}{coil} when $i = \pm I_0$.
\end{proposition}

\begin{proof}
Differentiate: $\frac{dE}{dt} = \frac{q}{C}\frac{dq}{dt} +
Li\frac{di}{dt} = i\left(\frac{q}{C} + L\frac{d^2q}{dt^2}\right) = 0$
by \cref{prop:g12:rlc-oscillations:equation}. Its value is read at
$t = 0$ (all \omterm{def:g12:rc-rl-circuits:capacitor}{capacitor}), and again at the all-current instant
$q = 0$ --- whence $\tfrac12 L I_0^2 = Q_0^2/2C$, recovering
$I_0 = Q_0/\sqrt{LC}$ by \omterm{def:g10:energy-conservation:energy}{energy} alone.
\end{proof}

\begin{omfigure}
\begin{tikzpicture}
\begin{axis}[omaxis, width=8.6cm, height=4.6cm, xlabel={$t/T_0$},
    ylabel={energy $/\,E$}, xmin=0, xmax=1, ymin=0, ymax=1.3,
    legend pos=south east]
  \addplot[omDef, thick, domain=0:1, samples=120] {cos(360*x)^2};
  \addplot[omProp, thick, domain=0:1, samples=120] {sin(360*x)^2};
  \addplot[omThm, dashed, thick, domain=0:1, samples=2] {1};
  \legend{$E_C$, $E_L$, $E$}
\end{axis}
\end{tikzpicture}

{\small One \omterm{def:g12:oscillators-and-time:oscillator}{period} of the waltz: \omterm{def:g10:energy-conservation:energy}{energy} sloshes from plates to \omterm{def:g11:magnetic-fields:solenoid}{coil}
and back twice per \omterm{def:g12:oscillators-and-time:oscillator}{period}, the sum pinned at $Q_0^2/2C$ --- the very
figure of the \omterm{prop:g12:oscillators-and-time:spring-period}{spring--mass oscillator}, relabeled.}
\end{omfigure}

\begin{example}[The waltz, measured]\label{ex:g12:rlc-oscillations:waltz}
$C = \qty{1.0}{\micro F}$ charged to $U_0 = \qty{10}{V}$, switched onto
$L = \qty{10}{mH}$: \omterm{def:g10:energy-conservation:energy}{energy} $E = \tfrac12 C U_0^2 = \qty{5.0e-5}{J}$,
\omterm{def:g12:oscillators-and-time:oscillator}{period} $T_0 = \qty{0.63}{ms}$, and a peak current
$I_0 = U_0\sqrt{C/L} = \qty{0.10}{A}$ reached a quarter \omterm{def:g12:oscillators-and-time:oscillator}{period}
(\qty{0.16}{ms}) after the switch closes.
\end{example}

\section{Real circuits: damped oscillations}

Every real \omterm{def:g11:magnetic-fields:solenoid}{coil} is wound from \omterm{def:g11:work-of-force:sign}{resistive} wire: the honest loop is a
series \emph{RLC} circuit, and the \omterm{def:g11:circuits-and-power:resistance}{resistance} taxes the waltz.

\begin{definition}[Pseudo-periodic and aperiodic regimes]\label{def:g12:rlc-oscillations:regimes}
\index{RLC circuit}\index{pseudo-periodic regime}\index{aperiodic regime}
In a series RLC circuit the resistor dissipates \omterm{def:g10:energy-conservation:energy}{energy} by Joule
heating (\cref{ch:g11:circuits-and-power}): the \omterm{def:g12:oscillators-and-time:oscillator}{free oscillations} are
\emph{\omterm{def:g12:oscillators-and-time:damping}{damped}}. Two regimes, read off an \omterm{def:g10:signals-and-waves:oscilloscope}{oscilloscope}:
\begin{itemize}
  \item \emph{\omterm{def:g10:signals-and-waves:period}{pseudo-periodic}} (small $R$): the circuit still
        oscillates inside a shrinking envelope, with a repeat time ---
        the \emph{pseudo-period}\index{pseudo-period!electrical} ---
        close to $T_0 = 2\pi\sqrt{LC}$ when the \omterm{def:g12:oscillators-and-time:damping}{damping} is light;
  \item \emph{aperiodic} (large $R$): the charge creeps back to zero
        without ever overshooting --- no oscillation at all.
\end{itemize}
The exact mirror of the \omterm{def:g12:oscillators-and-time:damping}{damped} \omterm{def:g12:oscillators-and-time:oscillator}{mechanical oscillator}
(\cref{def:g12:oscillators-and-time:damping}).
\end{definition}

\begin{omfigure}
\begin{tikzpicture}
\begin{axis}[omaxis, width=9.2cm, height=5.2cm,
    xlabel={$t$ (\unit{\micro s})}, ylabel={$u$ (\unit{V})},
    xmin=0, xmax=10, ymin=-6.6, ymax=8.2, legend pos=north east]
  \addplot[omDef, thick, domain=0:10, samples=300]
    {6*exp(-x/3)*cos(deg(2*pi*x/1.6))};
  \addplot[omProp, thick, domain=0:10, samples=100] {6*exp(-x/1.2)};
  \legend{pseudo-periodic (small $R$), aperiodic (large $R$)}
  \addplot[omThm, dashed, domain=0:10, samples=80, forget plot]
    {6*exp(-x/3)};
  \addplot[omThm, dashed, domain=0:10, samples=80, forget plot]
    {-6*exp(-x/3)};
\end{axis}
\end{tikzpicture}

{\small The two fates of a real \omterm{def:g12:rlc-oscillations:regimes}{RLC circuit}: oscillations dying inside
an exponential envelope (dashed), or a creep to zero with no
oscillation. Lightly \omterm{def:g12:oscillators-and-time:damping}{damped}, the \omterm{def:g12:rlc-oscillations:regimes}{pseudo-period} stays $\approx T_0$.}
\end{omfigure}

\begin{remark}[Maintaining oscillations]\label{rem:g12:rlc-oscillations:maintain}
\index{maintained oscillations}
To keep the \omterm{def:g12:oscillators-and-time:oscillator}{amplitude} alive, feed the circuit exactly the \omterm{def:g10:energy-conservation:energy}{energy} the
\omterm{def:g11:circuits-and-power:resistance}{resistance} burns, once per cycle and \emph{in step} with the swing: an
amplifier does for the RLC loop what the escapement's push does for
the pendulum (\cref{ch:g12:oscillators-and-time}). Every quartz watch
and every radio transmitter is such a maintained \omterm{def:g12:oscillators-and-time:oscillator}{oscillator}, running
at its \omterm{def:g12:rlc-oscillations:period}{natural frequency} $f_0$.
\end{remark}

\section{The mechanical twin}

\begin{proposition}[The electrical--mechanical dictionary]\label{prop:g12:rlc-oscillations:analogy}
\index{electrical--mechanical analogy}
The LC equation is the spring--mass equation
(\cref{prop:g12:oscillators-and-time:spring-period}) word for word:
\begin{center}\small
\begin{tabular}{lcl}
\omterm{prop:g12:rlc-oscillations:equation}{LC circuit} & & spring--mass glider\\[2pt]
charge $q$ & $\longleftrightarrow$ & position $x$\\
current $i = dq/dt$ & $\longleftrightarrow$ & velocity $v = dx/dt$\\
\omterm{def:g12:rc-rl-circuits:inductor}{inductance} $L$ & $\longleftrightarrow$ & mass $m$\\
inverse \omterm{def:g12:rc-rl-circuits:capacitor}{capacitance} $1/C$ & $\longleftrightarrow$ & \omterm{def:g12:oscillators-and-time:spring}{stiffness} $k$\\
\omterm{def:g11:circuits-and-power:resistance}{resistance} $R$ & $\longleftrightarrow$ & \omterm{def:g10:inertia:inventory}{friction}\\
$L\,\frac{d^2q}{dt^2} + \frac{q}{C} = 0$ & $\longleftrightarrow$ &
  $m\,\frac{d^2x}{dt^2} + kx = 0$\\[2pt]
$T_0 = 2\pi\sqrt{LC}$ & $\longleftrightarrow$ &
  $T_0 = 2\pi\sqrt{m/k}$\\[2pt]
$E_C = q^2/2C$ & $\longleftrightarrow$ & $E_p = \tfrac12 kx^2$\\
$E_L = \tfrac12 Li^2$ & $\longleftrightarrow$ & $E_k = \tfrac12 mv^2$
\end{tabular}
\end{center}
Same equation, same solutions: every result about either \omterm{def:g12:oscillators-and-time:oscillator}{oscillator}
translates instantly into the other language.
\end{proposition}

\begin{proof}
Compare \cref{prop:g12:rlc-oscillations:equation} with
$m\,d^2x/dt^2 = -kx$: substituting $q \to x$, $L \to m$, $1/C \to k$
turns one into the other, and $LC \to m/k$ turns
$2\pi\sqrt{LC}$ into $2\pi\sqrt{m/k}$.
\end{proof}

\begin{example}[Tuning a radio]\label{ex:g12:rlc-oscillations:tuning}
An aerial feeds a whisper of \emph{every} station into an LC loop, but
only the one broadcasting near $f_0$ pushes in step and builds up ---
\omterm{def:g12:oscillators-and-time:resonance}{resonance} (\cref{def:g12:oscillators-and-time:resonance}). The dial
rotates the plates of a variable \omterm{def:g12:rc-rl-circuits:capacitor}{capacitor}: changing $C$ moves
$f_0 = 1/(2\pi\sqrt{LC})$, and the needle hands the earphone one
station at a time. On the transmitting side, a maintained \omterm{def:g12:oscillators-and-time:oscillator}{oscillator}
(\cref{rem:g12:rlc-oscillations:maintain}) sets the broadcast
\omterm{def:g12:oscillators-and-time:oscillator}{frequency}; in a watch, a quartz crystal --- a mechanical twin of
breathtaking \omterm{def:g12:oscillators-and-time:spring}{stiffness} --- plays the same role.
\end{example}

\section{Exercises}

\begin{exercise}[$\star$]\label{exo:g12:rlc-oscillations:1}
Compute $T_0$ and $f_0$ for (a) $L = \qty{10}{mH}$,
$C = \qty{1.0}{\micro F}$; (b) $L = \qty{1.0}{mH}$,
$C = \qty{100}{pF}$. By what factor do the frequencies differ?
\end{exercise}

\begin{exercise}[$\star$]\label{exo:g12:rlc-oscillations:2}
Using $\qty{1}{H} = \qty{1}{V.s/A}$ and $\qty{1}{F} = \qty{1}{C/V}$,
show \omterm{def:g10:orders-of-magnitude:unit}{unit} by \omterm{def:g10:orders-of-magnitude:unit}{unit} that $\sqrt{LC}$ is a time. Why does the $2\pi$
change nothing?
\end{exercise}

\begin{exercise}[$\star$]\label{exo:g12:rlc-oscillations:3}
A \qty{470}{nF} \omterm{def:g12:rc-rl-circuits:capacitor}{capacitor} is charged to \qty{6.0}{V}. Compute $Q_0$
and the stored \omterm{def:g10:energy-conservation:energy}{energy}; describe, without computation, what follows
when it is switched onto an ideal \omterm{def:g11:magnetic-fields:solenoid}{coil}.
\end{exercise}

\begin{exercise}[$\star$]\label{exo:g12:rlc-oscillations:4}
An \omterm{prop:g12:rlc-oscillations:equation}{LC circuit} follows $q(t) = 2.0\cos(2\pi t/T_0)$ in microcoulombs,
with $T_0 = \qty{4.0}{ms}$. Read off $Q_0$; compute $f_0$, $q(T_0/2)$,
$i(0)$, and the peak current $I_0 = 2\pi Q_0/T_0$.
\end{exercise}

\begin{exercise}[$\star$]\label{exo:g12:rlc-oscillations:5}
At which instants of the cycle is the \omterm{def:g10:energy-conservation:energy}{energy} entirely in the
\omterm{def:g12:rc-rl-circuits:capacitor}{capacitor}? Entirely in the \omterm{def:g11:magnetic-fields:solenoid}{coil}? What is the mechanical counterpart of
each moment for a pendulum?
\end{exercise}

\begin{exercise}[$\star\star$]\label{exo:g12:rlc-oscillations:6}
Verify by substitution that $q(t) = Q_0\cos(2\pi t/T_0)$ satisfies
$L\,d^2q/dt^2 + q/C = 0$ exactly when $T_0 = 2\pi\sqrt{LC}$.
\end{exercise}

\begin{exercise}[$\star\star$]\label{exo:g12:rlc-oscillations:7}
A \omterm{def:g10:signals-and-waves:wave}{long-wave} station broadcasts at \qty{198}{kHz}. With
$C = \qty{220}{pF}$, what \omterm{def:g12:rc-rl-circuits:inductor}{inductance} tunes to it?
\end{exercise}

\begin{exercise}[$\star\star$]\label{exo:g12:rlc-oscillations:8}
$C = \qty{1.0}{\micro F}$ charged to \qty{10}{V} discharges into
$L = \qty{10}{mH}$. Compute the stored \omterm{def:g10:energy-conservation:energy}{energy}, the peak current (by
\omterm{def:g10:energy-conservation:energy}{energy} conservation), and the instant it is first reached.
\end{exercise}

\begin{exercise}[$\star\star$]\label{exo:g12:rlc-oscillations:9}
Explain why $E_C$ and $E_L$ oscillate at \omterm{def:g12:oscillators-and-time:oscillator}{frequency} $2f_0$, not $f_0$.
Show that at $t = T_0/8$ (started from $q = Q_0$) the \omterm{def:g10:energy-conservation:energy}{energy} is split
exactly in half.
\end{exercise}

\begin{exercise}[$\star\star$]\label{exo:g12:rlc-oscillations:10}
On an \omterm{def:g10:signals-and-waves:oscilloscope}{oscillogram} of a lightly \omterm{def:g12:oscillators-and-time:damping}{damped} \omterm{def:g12:rlc-oscillations:regimes}{RLC circuit}, successive maxima
read \qty{4.0}{V} then \qty{3.0}{V}. Assuming the ratio stays
constant, predict the sixth maximum after the \qty{4.0}{V} one. What
fraction of the \emph{\omterm{def:g10:energy-conservation:energy}{energy}} survives each \omterm{def:g12:rlc-oscillations:regimes}{pseudo-period}? What can
you say of the \omterm{def:g12:rlc-oscillations:regimes}{pseudo-period} itself?
\end{exercise}

\begin{exercise}[$\star\star$]\label{exo:g12:rlc-oscillations:11}
An \omterm{prop:g12:rlc-oscillations:equation}{LC circuit} has $L = \qty{0.50}{H}$ and $C = \qty{20}{\micro F}$.
Compute $T_0$; then, using the dictionary
(\cref{prop:g12:rlc-oscillations:analogy}), find the \omterm{def:g12:oscillators-and-time:spring}{stiffness} $k$
giving a \qty{0.50}{kg} glider the same \omterm{def:g12:oscillators-and-time:oscillator}{period}, and check it.
\end{exercise}

\begin{exercise}[$\star\star\star$]\label{exo:g12:rlc-oscillations:12}
A tuner uses $L = \qty{0.20}{mH}$ and a variable \omterm{def:g12:rc-rl-circuits:capacitor}{capacitor} sweeping
\qtyrange{50}{500}{pF}. Compute the two extreme frequencies. Show that
$f_0 \propto 1/\sqrt{C}$, so tenfold in $C$ gives only
$\sqrt{10} \approx 3.2$ in \omterm{def:g12:oscillators-and-time:oscillator}{frequency}.
\end{exercise}

\begin{exercise}[$\star\star\star$]\label{exo:g12:rlc-oscillations:13}
A real \omterm{def:g11:magnetic-fields:solenoid}{coil} makes the circuit lose $5.0\%$ of its \omterm{def:g10:energy-conservation:energy}{energy} each
\omterm{def:g12:rlc-oscillations:regimes}{pseudo-period}. After how many \omterm{def:g12:oscillators-and-time:oscillator}{periods} has the \omterm{def:g10:energy-conservation:energy}{energy} halved? What is
the \omterm{def:g12:oscillators-and-time:oscillator}{amplitude} then, as a fraction of the start? At
$f_0 = \qty{1.0}{MHz}$, how long is that in microseconds?
\end{exercise}

\begin{exercise}[$\star\star\star$]\label{exo:g12:rlc-oscillations:14}
A watch \omterm{def:g12:oscillators-and-time:oscillator}{oscillator} stores $E = \qty{1.0e-9}{J}$ and loses $2.0\%$ of
it per cycle at $f_0 = \qty{32768}{Hz}$. What average \omterm{def:g11:work-of-force:power}{power} must the
maintaining circuit supply? Why does the push have to arrive in step
with the oscillation?
\end{exercise}

\begin{exercise}[$\star\star\star$]\label{exo:g12:rlc-oscillations:15}
After each edge of a square \omterm{def:g10:signals-and-waves:signal}{signal}, an experimenter's wiring ``rings''
at \qty{25}{MHz}; the stray \omterm{def:g12:rc-rl-circuits:inductor}{inductance} is about \qty{1.0}{\micro H}.
Estimate the stray \omterm{def:g12:rc-rl-circuits:capacitor}{capacitance} responsible. Which regime should a
small added series resistor push the circuit into, and why does that
cure the ringing?
\end{exercise}

\section{Problem: Building a crystal radio}

\begin{problem}[{Weekend problem --- building a crystal radio: one
coil, one variable capacitor, a diode and an earphone --- no battery,
no amplifier, and the evening news arrives on the energy of the wave
itself}]
\label{pb:g12:rlc-oscillations:1}
A kit from a grandparent's attic: a \omterm{def:g11:magnetic-fields:solenoid}{coil} $L = \qty{0.20}{mH}$ wound on
a cardboard tube, a variable \omterm{def:g12:rc-rl-circuits:capacitor}{capacitor} whose plates rotate from
\qty{20}{pF} to \qty{480}{pF}, a crystal diode, an earphone, and
\qty{20}{m} of aerial wire. The stations to catch broadcast between
\qty{530}{kHz} and \qty{1700}{kHz}.

\medskip
\noindent\textbf{Part I --- The tuning loop.}
\begin{enumerate}
  \item The \omterm{def:g12:rc-rl-circuits:capacitor}{capacitor}, charged, is left to discharge into the \omterm{def:g11:magnetic-fields:solenoid}{coil}:
        write the loop law and turn it into the differential equation
        for $q(t)$.
  \item Verify by substitution that $q(t) = Q_0\cos(2\pi t/T_0)$
        solves it, provided $T_0 = 2\pi\sqrt{LC}$.
  \item Check by dimensional analysis that $\sqrt{LC}$ is a time.
  \item Dial set to $C = \qty{330}{pF}$: compute $T_0$ and $f_0$. Is
        the loop tuned inside the band?
  \item Find the \omterm{def:g12:rc-rl-circuits:capacitor}{capacitance} that tunes to exactly \qty{1.00}{MHz}.
\end{enumerate}

\medskip
\noindent\textbf{Part II --- The variable \omterm{def:g12:rc-rl-circuits:capacitor}{capacitor}.}
\begin{enumerate}[resume]
  \item Compute the \omterm{def:g12:rc-rl-circuits:capacitor}{capacitance} needed for \qty{530}{kHz}.
  \item Compute the \omterm{def:g12:rc-rl-circuits:capacitor}{capacitance} needed for \qty{1700}{kHz}.
  \item Compute the two frequencies the kit's \qtyrange{20}{480}{pF}
        \omterm{def:g12:rc-rl-circuits:capacitor}{capacitor} actually reaches. Does it cover the whole band?
  \item Show that $f_0 \propto 1/\sqrt{C}$ at fixed $L$; by what
        factor must $C$ change to double $f_0$?
  \item The plates open linearly, so $C$ falls steadily as the dial
        turns. Explain why the stations then crowd together at one end
        of the dial --- and which end.
\end{enumerate}

\medskip
\noindent\textbf{Part III --- \omterm{def:g10:energy-conservation:energy}{Energy} and \omterm{def:g12:oscillators-and-time:damping}{damping}.}
Tuned to \qty{1.00}{MHz} ($C = \qty{127}{pF}$), the aerial momentarily
charges the \omterm{def:g12:rc-rl-circuits:capacitor}{capacitor} to $U_0 = \qty{20}{mV}$.
\begin{enumerate}[resume]
  \item Compute the charge $Q_0$ and the stored \omterm{def:g10:energy-conservation:energy}{energy} $E$.
  \item Compute the peak current $I_0$, using \omterm{def:g10:energy-conservation:energy}{energy} conservation.
  \item The \omterm{def:g11:magnetic-fields:solenoid}{coil}'s \omterm{def:g11:circuits-and-power:resistance}{resistance} dissipates $4.0\%$ of the stored \omterm{def:g10:energy-conservation:energy}{energy}
        per cycle: after how many cycles has the \omterm{def:g10:energy-conservation:energy}{energy} halved?
  \item How long is that in microseconds? Why can this ``ringing''
        alone not play the evening news?
  \item Explain why the arriving \omterm{def:g10:signals-and-waves:wave}{wave} keeps the loop swinging only
        when tuned to that station --- and what maintains the
        oscillation at the \emph{transmitter}
        (\cref{rem:g12:rlc-oscillations:maintain}).
\end{enumerate}

\medskip
\noindent\textbf{Part IV --- The mechanical twin.}
\begin{enumerate}[resume]
  \item Using the dictionary (\cref{prop:g12:rlc-oscillations:analogy})
        with a glider of mass $m = \qty{0.20}{kg}$, compute the
        \omterm{def:g12:oscillators-and-time:spring}{stiffness} $k$ of the spring matching the tuned circuit.
  \item A stiff car spring has $k \approx \qty{5e4}{N/m}$. Comment:
        can a bench-top spring--mass system oscillate at megahertz?
  \item Keep a reasonable $k = \qty{100}{N/m}$ instead: what mass
        oscillates at \qty{1.00}{MHz}?
  \item What real object is such a featherweight, ultra-stiff
        \omterm{def:g12:oscillators-and-time:oscillator}{mechanical oscillator}, and where did you meet it earlier this
        year (\cref{ch:g12:oscillators-and-time})?
  \item Spec sheet, one sentence: the band your radio sweeps, the
        \omterm{def:g12:rc-rl-circuits:capacitor}{capacitance} that picks \qty{1.00}{MHz}, how long a lone ring
        lasts, and what plays the news anyway.
\end{enumerate}
\end{problem}
